\documentclass[11pt,a4paper]{article}
\usepackage[utf8]{inputenc}
\usepackage[spanish]{babel}
\usepackage{amsmath,amsfonts,amssymb}
\usepackage{graphicx}
\usepackage{geometry}
\geometry{margin=2.5cm}

\title{\textbf{Guía 0: Diseño de un nodo micro-ROS para motor DC}\\Control Inteligente -- ING01343-ING278}
\author{Politécnico Colombiano Jaime Isaza Cadavid\\Profesor: Deimer Miranda Montoya, MSc.(c).}
\date{}

\begin{document}
\maketitle

\section{Introducción}
En esta práctica se diseña e implementa la arquitectura conceptual de un nodo micro-ROS sobre un ESP32 para interactuar con un motor DC Pololu, utilizando un puente H L298 y un encoder incremental de cuadratura.

\section{Arquitectura del Sistema}
\begin{itemize}
    \item \textbf{ESP32 (micro-ROS):} nodo \texttt{motor\_step\_node}, recepción de \texttt{/pwm\_input}, cálculo de velocidades y publicación en \texttt{/vel\_rad\_s} y \texttt{/vel\_rpm}.
    \item \textbf{PC (ROS 2):} micro\_ros\_agent, nodos de adquisición (\texttt{data\_base\_node}) y visualización (\texttt{graph\_vel\_node}).
\end{itemize}

\section{Contrato de Interfaces}
\begin{table}[h!]
\centering
\begin{tabular}{|l|l|l|l|}
\hline
\textbf{Tópico} & \textbf{Dirección} & \textbf{Tipo} & \textbf{Unidad} \\ \hline
\texttt{/pwm\_input} & Entrada (PC $\to$ ESP32) & \texttt{std\_msgs/msg/Float32} & \% ($-100$ a $100$) \\ \hline
\texttt{/vel\_rad\_s} & Salida (ESP32 $\to$ PC) & \texttt{std\_msgs/msg/Float32} & rad/s \\ \hline
\texttt{/vel\_rpm} & Salida (ESP32 $\to$ PC) & \texttt{std\_msgs/msg/Float32} & rpm \\ \hline
\end{tabular}
\caption{Interfaces ROS 2 del sistema.}
\end{table}

\end{document}
